\section{Marco Teórico}

El análisis de señales EEG constituye una herramienta en el estudio de la actividad cerebral y en el apoyo al diagnóstico de trastornos neurológicos y mentales, debido a su carácter no invasivo, alta resolución temporal y bajo costo relativo \cite{Nunez2006,Acharya2016}. Estas señales reflejan la dinámica eléctrica del cerebro y se caracterizan por presentar comportamientos altamente no lineales, no estacionarios y dependientes del contexto, lo que dificulta su modelado mediante enfoques clásicos de análisis de señales.
En el ámbito clínico, múltiples estudios han identificado patrones diferenciales en señales EEG asociados a diversos trastornos mentales, incluyendo TDAH, depresión, esquizofrenia y trastornos del espectro autista. Estos hallazgos han motivado el uso de técnicas de aprendizaje automático para el análisis automatizado de EEG, permitiendo la extracción de características discriminativas sin depender exclusivamente de conocimiento experto \cite{Craik2019,Bashivan2016}. Sin embargo, estos enfoques suelen requerir grandes volúmenes de datos etiquetados y presentan limitaciones en términos de generalización entre sujetos, sesiones y condiciones de adquisición.
Para abordar estas limitaciones, han surgido enfoques basados en aprendizaje profundo, los cuales permiten modelar representaciones jerárquicas directamente a partir de los datos. No obstante, estos modelos son altamente dependientes de datos anotados y tienden a sobreajustarse a distribuciones específicas, lo que limita su aplicabilidad en escenarios reales donde existe variabilidad significativa. En este contexto, el aprendizaje auto-supervisado ha emergido como una estrategia efectiva para aprovechar grandes volúmenes de datos no etiquetados, permitiendo el aprendizaje de representaciones latentes transferibles a múltiples tareas \cite{Mathew2024,Abbaspourazad2023}.
De manera complementaria, el aprendizaje por transferencia (transfer learning) ha sido ampliamente utilizado para mejorar el desempeño de modelos en escenarios con datos limitados, reutilizando conocimiento adquirido en dominios relacionados \cite{Pan2010}. En el análisis de señales EEG, este enfoque permite adaptar modelos entrenados en grandes bases de datos a nuevas poblaciones o condiciones clínicas. Sin embargo, su efectividad depende de la similitud entre dominios y puede verse afectada por el fenómeno de cambio de distribución (dataset shift), donde las diferencias en los datos de entrenamiento y prueba degradan el desempeño del modelo \cite{MorenoTorres2012}.
En este sentido, la generalización de dominio y la adaptación de dominio han sido propuestas como estrategias para mejorar la robustez de los modelos frente a la variabilidad en los datos \cite{Wang2022DG}. Estas técnicas buscan aprender representaciones invariantes a las condiciones de adquisición, sujetos o dispositivos, lo cual resulta relevante en señales EEG, donde la heterogeneidad es inherente.
Por otra parte, un aspecto crítico en aplicaciones clínicas es la confiabilidad de las predicciones generadas por los modelos. En este contexto, la cuantificación de incertidumbre se ha convertido en un componente clave para el desarrollo de sistemas de apoyo a la decisión. Diversos enfoques han sido propuestos, incluyendo métodos bayesianos, aproximaciones como Monte Carlo Dropout, técnicas basadas en ensamblados y métodos de calibración como la predicción conformal \cite{Loftus2022,ovadia2019trustmodelsuncertaintyevaluating}. Estos métodos permiten estimar la incertidumbre asociada a las predicciones, lo que permite asignar un nivel de confianza en el pronóstico del modelo.
Adicionalmente, los procesos gaussianos representan un enfoque no paramétrico dentro del marco bayesiano que permite modelar distribuciones sobre funciones, proporcionando estimaciones probabilísticas con incertidumbre explícita \cite{Rasmussen2006}. No obstante, su aplicación en escenarios de alta dimensionalidad y grandes volúmenes de datos puede resultar limitada en términos de escalabilidad, lo que ha motivado el desarrollo de aproximaciones híbridas y alternativas más eficientes.
A pesar de estos avances, persisten desafíos importantes en la integración de estas técnicas dentro de un marco unificado. En particular, la combinación de aprendizaje auto-supervisado, transferencia de conocimiento, generalización entre dominios y estimación de incertidumbre en el análisis de señales EEG sigue siendo un área abierta de investigación \cite{ABDAR2021243} y actualmente no existe un método generalizable basado en EEG para apoyo diagnóstico clínico. Esta brecha es especialmente relevante en contextos clínicos, donde la variabilidad de los datos y la necesidad de interpretabilidad y confiabilidad limitan la adopción de sistemas basados en inteligencia artificial.